\section{Introdução e Objetivos}

\begin{frame}

    \frametitle{Introdução}

    \begin{itemize}
	 	\item A classificação manual de resíduos é uma atividade que demanda tempo e pode expor trabalhadores a condições insalubres.
	 	\item Redes neurais convolucionais permitem automatizar a classificação de imagens.
    	\item O projeto buscou desenvolver uma solução que considerasse não apenas o desempenho, mas também a generalização e a explicabilidade das classificações.    	
    \end{itemize} 
    
\end{frame}

\begin{frame}

    \frametitle{Objetivos}

    \begin{itemize}
    
        \item Objetivo Geral
            \begin{itemize}
                \item Desenvolver e avaliar um classificador de resíduos baseado em redes neurais convolucionais que fosse generalizável entre diferentes fontes de imagens e interpretável.
            \end{itemize}
            
        \item Objetivos Específicos
            \begin{itemize}
            	\item Construir um conjunto de dados integrado, com rótulos harmonizados e sem vazamento entre treino, validação e teste.         	
            	\item Desenvolver um protocolo de avaliação entre diferentes domínios.
            	\item Comparar arquiteturas convolucionais.
            	\item Investigar as decisões dos modelos utilizando Grad-CAM.
            	\item Desenvolver uma aplicação móvel para utilização do modelo.
                                
            \end{itemize}
            
    \end{itemize}

\end{frame}